\section{Introduction}
\label{sec:intro}

\IEEEPARstart{P}{reemptible} capacity is the cheapest way to rent an accelerator and the least dependable. A spot
instance is reclaimed with a warning of two minutes on one large cloud and at most thirty seconds on two
others~\cite{wu2024cantbelate,lee2022spotlake}. In that window a training job has to get its state off the machine. For
a model trained with Adam~\cite{kingma2015adam,loshchilov2019adamw} the state is the weights and two moment tensors of
the same shape, twelve bytes per parameter in single precision. A \llamaParams{}-billion-parameter model therefore owns
\llamaFullGB{}\,GB of state, which takes \llamaFullDrain{}\,s to drain over a 10\,Gb/s link before any serialisation
cost: inside a two-minute notice, outside a thirty-second one.

The systems that keep training alive on such machines differ in almost everything except this. They re-plan pipelines
when workers vanish~\cite{athlur2022varuna,thorpe2023bamboo,jang2023oobleck,duan2024parcae}, replicate state to
peers~\cite{wang2023gemini}, rebuild it from mirrored gradients~\cite{bhardwaj2025checkmate}, migrate it directly
between nodes~\cite{liu2026etc}, or snapshot the whole device~\cite{stoyanov2025criugpu,wei2025phoenixos}. All of them
move the weights and both moments in full. The most recent migration system states the assumption outright: ``it is
necessary to migrate model parameters, optimizer states and other states''~\cite{liu2026etc}. The schedulers that decide
when to use spot capacity then fold everything that happens at a switch into one constant, the changeover
delay~\cite{wu2024cantbelate}.

Work on optimizers points the other way. Adafactor trains with only the row and column sums of the second
moment~\cite{shazeer2018adafactor}; later work removes or compresses most of it for the whole
run~\cite{zhang2025adammini,kalra2025slimadam} and stores both moments in eight or four
bits~\cite{dettmers2022eightbit,li2023fourbit}. If a run can live with a coarse second moment at every step, it can
presumably survive receiving a coarse one once. Nobody has measured that. The question is not whether the state can be
made smaller, which checkpoint compression already answers~\cite{li2024excp,zhang2026guarded}, but which of it has to
exist on the wire at all, and what is paid in training progress for each part left behind.

This paper answers that question by experiment. A run is stopped at step $T$, its state is reduced according to a
\emph{strategy}, and training resumes from what is left. The resumed run is compared, step by step, with a control that
was never interrupted. Because data order, dropout masks and the learning-rate schedule are all pure functions of the
step index, the control and the resumed run differ in exactly one thing, the state that crossed, and a resume from the
full state reproduces the control bit for bit. Because two runs that differ only in their seed also drift apart, every
gap is read against that spread, the \emph{seed band}; a migrated run that stays inside it cannot be told from a
different seed of the same job.

The contributions are as follows.

\begin{itemize}
\item \textbf{A controlled ablation of optimizer state at a migration boundary.} \nArms{} payload strategies, from the
full state down to the weights alone, are evaluated over \nCells{} migrations on three model sizes, four points in
training and up to \nSeedsMax{} seeds, with bootstrap intervals over seeds and censored recovery times
(Sections~\ref{sec:setup} and~\ref{sec:results}).
\item \textbf{An account of the arrival transient.} With both moments zeroed, the first update is not small, as is
sometimes assumed, but about $\sqrt{(1-\beta_1^2)/(1-\beta_2)}$ times the steady one, and it decays on the horizon of
the second moment if the step counter is kept and of the first if it is reset. The closed forms agree with measured
update sizes, and explain why the first moment may be dropped, why the second may be sketched, and why a first moment
re-estimated at the destination is harmful (Section~\ref{sec:transient}).
\item \textbf{Measured bytes and latency.} Payload sizes are taken from the serialised bundle and agree with the
analytic count to within \bytesWorstPct\% ; drain and restore are timed per stage over repeated rounds
(Section~\ref{sec:results}).
\item \textbf{Migration across device classes.} The same payload is resumed on a CPU, a discrete GPU and an integrated
GPU of a different vendor, and the gap between devices is itself measured against the seed band
(Section~\ref{sec:device}).
\item \textbf{A policy.} A cost model takes the measured curve, a notice window, a bandwidth and a price, and returns the
cheapest strategy that fits (Section~\ref{sec:policy}). Its outputs for models larger than those trained here are
extrapolations and are labelled as such.
\end{itemize}

The experiments run on one laptop with models of \paramsXS{} to \paramsM{} million parameters, and that limit is
stated wherever it matters. What the paper offers at that scale is a method, a noise floor and a set of mechanisms that
make predictions; Section~\ref{sec:limits} lists what remains to be checked on a real cluster.
